% B3: organização visual do autor; não reproduz figura dos artigos.
\begin{tikzpicture}[
    etapa/.style={draw=black!65, rounded corners=2pt, align=center,
        text width=3.9cm, minimum height=2.0cm, inner sep=7pt,
        font=\small,fill=blue!4}, node distance=8mm,
    seta/.style={-{Latex[length=2.5mm]},thick}]
\node[etapa] (ler) {\textbf{Ler}\\[5pt]Identificar o conteúdo da mensagem.};
\node[etapa,right=of ler] (entender)
    {\textbf{Compreender}\\[5pt]Relacionar a informação ao problema no programa.};
\node[etapa,right=of entender] (formular)
    {\textbf{Formular uma ação}\\[5pt]Planejar uma alteração coerente com o diagnóstico.};
\draw[seta] (ler) -- (entender);
\draw[seta] (entender) -- (formular);
\node[below=6mm of entender,align=center,text width=12.5cm,font=\small]
    {As etapas organizam a análise de respostas aos erros;\newline
    o percurso não representa uma taxa de sucesso nem uma medição do Luminthor.};
\end{tikzpicture}
